\section{Combinatorics and Discrete Mathematics}

\begin{notationtable}
$[n]=\{1,2,\dots,n\}$ & Standard $n$-element set, especially in combinatorics; $[0]=\varnothing$ when needed. \\
$\displaystyle\binom{n}{r},\ {}_nC_r$ & Number of $r$-element subsets of an $n$-element set. \\*
$\displaystyle\left(\!\binom{n}{r}\!\right),\ {}_nH_r$ & $r$-combinations with repetition from $n$ types; equals $\displaystyle\binom{n+r-1}{r}$. \\
$\displaystyle\binom{n}{n_1,\dots,n_k}$ & Multinomial coefficient, for nonnegative integers $n_1,\dots,n_k$ with $n_1+\cdots+n_k=n$. \\
$n!$ & Factorial. \\
$n^{\underline{k}}$ & Falling factorial, $n(n-1)\cdots(n-k+1)$. \\*
$n^{\overline{k}}$ & Rising factorial, $n(n+1)\cdots(n+k-1)$. \\
$(a)_k$ & Rising Pochhammer symbol $a(a+1)\cdots(a+k-1)$ in special-functions contexts; it is not used here for the falling factorial. \\
\addlinespace[0.35em]
$(X,\preceq)$ & Partially ordered set (poset). \\
$\prec,\ \succeq,\ \succ$ & Strict and reverse order relations associated with $\preceq$. \\
$\le_{\mathrm{lex}}$ & Lexicographic order. \\
$x\lessdot y$ & Cover relation: $y$ covers $x$. \\
$B_n$ & Boolean lattice (Boolean poset) on an $n$-element ground set. \\
\addlinespace[0.35em]
$G=(V,E)$ & Graph with vertex set $V$ and edge set $E$. \\
$V(G)$ & Vertex set of $G$. \\
$E(G)$ & Edge set of $G$. \\
$N(v),\ \Gamma(v)$ & Neighborhood of vertex $v$. \\
$\deg_G(v),\ \deg(v)$ & Degree of vertex $v$; the graph subscript may be omitted when clear. \\
$\deg^+(v),\ \deg^-(v)$ & Out-degree and in-degree of $v$ in a directed graph. \\
$d_G(u,v),\ d(u,v)$ & Graph distance between vertices $u$ and $v$. \\
$G[S]$ & Subgraph of $G$ induced by the vertex set $S\subseteq V(G)$. \\
$G-v,\ G-S$ & Graph obtained by deleting the vertex $v$, or all vertices in $S$, together with their incident edges. \\
$G-e$ & Graph obtained by deleting the edge $e$. \\
$G/e$ & Graph obtained by contracting the edge $e$. \\
$\overline G$ & Complement of the graph $G$. \\
$K_n$ & Complete graph on $n$ vertices. \\
$K_{m,n}$ & Complete bipartite graph with parts of sizes $m$ and $n$. \\
$P_n$ & Path graph on $n$ vertices. \\
$C_n$ & Cycle graph on $n$ vertices. \\
$\mathbf A(G),\ \mathbf A_G$ & Adjacency matrix of $G$. \\
$\mathbf D(G),\ \mathbf D_G$ & Degree matrix of $G$, the diagonal matrix of vertex degrees in the chosen vertex order. \\
$\mathbf L(G),\ \mathbf L_G$ & Graph Laplacian, usually $\mathbf D(G)-\mathbf A(G)$. \\
$\chi(G)$ & Chromatic number of $G$. \\
$\tau(G),\ T(G)$ & Number of spanning trees of $G$; $\tau(G)$ is preferred. \\
\addlinespace[0.35em]
$\begin{aligned} &R(k,\ell),\ r(k,\ell),\\[-0.15em] &R(k),\ r(k) \end{aligned}$ & Ramsey numbers; uppercase $R$ is preferred, and $R(k)=R(k,k)$ in the diagonal case. \\
$\displaystyle A(x)=\sum_{n=0}^{\infty}a_nx^n$ & Ordinary generating function of the sequence $(a_n)_{n\ge0}$. \\
$\displaystyle\sum_{n=0}^{\infty}a_n\frac{x^n}{n!}$ & Exponential generating function of the sequence $(a_n)_{n\ge0}$; the function name is chosen locally when needed. \\
$[x^n]F(x)$ & Coefficient of $x^n$ in the formal power series or generating function $F(x)$. \\
\end{notationtable}
